\chapter{Number Theory}

%%
%% CRT
%%
\section{The Chinese Remainder Theorem}
\label{appendix:crt}

\begin{definition}[Chinese Remainder Theorem]
    \label{def:crt}
    Let $q_0, q_1, \dots, q_{k-1}$ be a set of $k$ pairwise coprime positive integers, meaning $\gcd(q_i, q_j) = 1$ for all $i \neq j$. For any given sequence of arbitrary integers $x_0, x_1, \dots, x_{k-1}$, there exists a unique solution $x$ modulo the product of these moduli that satisfies the following system of simultaneous congruences:
    \begin{align*}
        x &\equiv x_0 \pmod{q_0} \\
        \vdots \\
        x &\equiv x_{k-1} \pmod{q_{k-1}}
    \end{align*}
    \begin{gather}
        \nonumber
        \text{Let } Q = \prod_{i=0}^{k-1}q_i, \; \hat{q}_i = \frac{Q}{q_i} \text{ and } \hat{q}_i^{-1} \text{ be the inverse s.t. } \hat{q}_i^{-1}\hat{q}_i \equiv 1 \pmod{q_i}
        \\
        \text{Then \gls{crt} states } x = \sum_{i=0}^{k-1} x_i \cdot \hat{q}_i \cdot \hat{q}_i^{-1} \pmod{Q} 
        \\
        \label{eq:crt-gadget}
        \text{Or equivalently } x = \langle \{ x_i\}_{i=0}^{k-1}, \{\hat{q}_i\cdot\hat{q}_i^{-1} \}_{i=0}^{k-1} \rangle \pmod{Q}
    \end{gather}
\end{definition}

%%
%% NTT
%%
\section{The Number Theoretic Transform}
\label{appendix:ntt}
The efficiency of modern ring-based homomorphic encryption relies strictly on substituting high-complexity polynomial arithmetic with low-complexity scalar operations. This appendix formalizes the algebraic mechanics of the Number Theoretic Transform (NTT) over the negacyclic polynomial ring $R_{q_i} = \mathbb{Z}_{q_i}[X]/(X^N+1)$, which dictates the behavior of a single RNS tower.

\newpage
\begin{definition}[Principal Roots of Unity]
    Let $N$ be a power of two and $q$ be an $N$-transformable modulus such that $q \equiv 1 \pmod{2N}$. The algebraic field $\mathbb{Z}_q$ is guaranteed to contain a \emph{primitive $2N$-th root of unity}, denoted $\psi$, which strictly satisfies $\psi^{2N} \equiv 1 \pmod q$ and $\psi^k \not\equiv 1 \pmod q$ for all $0 < k < 2N$. Consequently, it holds that $\psi^N \equiv -1 \pmod q$.
\end{definition}

\begin{definition}[Forward Negacyclic NTT]
    Let $a \in R_q$ be a polynomial defined by its coefficient vector $(c_0, c_1, \dots, c_{N-1})$. The roots of the cyclotomic polynomial $X^N+1 = 0$ in $\mathbb{Z}_q$ are explicitly the odd powers of $\psi$: $\{\psi^1, \psi^3, \dots, \psi^{2N-1}\}$. 
    
    The \emph{Forward NTT} is the linear evaluation mapping of $a(X)$ exclusively at these specific roots. The transformed vector $\hat{a} = \text{NTT}(a)$ is defined component-wise for $j \in \{0, 1, \dots, N-1\}$ as:
    \begin{equation}
        \label{eq:forward-ntt}
        \hat{a}_j = a(\psi^{2j+1}) = \sum_{i=0}^{N-1} c_i \psi^{(2j+1)i} \pmod q
    \end{equation}
\end{definition}

\begin{definition}[Inverse Negacyclic NTT]
    Because the evaluation points are distinct roots of unity, the forward mapping is a strict algebraic isomorphism and is thereby unconditionally invertible. Given the transformed vector $\hat{a}$, the \emph{Inverse NTT} (INTT) reconstructs the original polynomial coefficients $c_i$ for $i \in \{0, 1, \dots, N-1\}$ via the orthogonal properties of these roots. The inverse mapping is formulated as:
    \begin{equation}
        \label{eq:inverse-ntt}
        c_i = \text{INTT}(\hat{a})_i = N^{-1} \psi^{-i} \sum_{j=0}^{N-1} \hat{a}_j \psi^{-2ij} \pmod q
    \end{equation}
    where $N^{-1}$ denotes the modular multiplicative inverse of $N$ in $\mathbb{Z}_q$.
\end{definition}

\begin{theorem}[Negacyclic Convolution Theorem]
    Let $a, b \in R_q$. The coefficient vector of their polynomial product $c = a \cdot b \pmod{X^N+1}$ is canonically isomorphic to the Hadamard (element-wise) product of their respective NTT forms. Specifically:
    \begin{equation}
        c = \text{INTT}(\text{NTT}(a) \odot \text{NTT}(b))
    \end{equation}
    
    By exploiting the symmetry of roots of unity via fast decimation algorithms, functionally identical to the Fast Fourier Transform utilizing the Cooley-Tukey butterfly structure, the computational evaluation of both the forward and inverse NTT is mathematically reduced from a standard complexity of $\mathcal{O}(N^2)$ to $\mathcal{O}(N \log N)$~\cite{liang:2022}. % This precise algorithmic reduction forms the absolute baseline requirement for practical latency in RNS-variant homomorphic cryptography.
\end{theorem}